\chapter{Set-Theoretic Foundations}

\section{$\pi$-Systems}

\begin{definition}[$\pi$-System]
A collection $\mathcal{P} \subseteq \Powerset{X}$ is a \textbf{$\pi$-system} if it is closed under finite intersections:
\[ A, B \in \mathcal{P} \implies A \cap B \in \mathcal{P} \]
\end{definition}

\begin{intuition}
A $\pi$-system is the minimal structure capturing ``intersection-closed'' collections. The name comes from ``product'' (intersection corresponds to logical AND). $\pi$-systems are important because:
\begin{itemize}
    \item They arise naturally: intervals, rectangles, open sets form $\pi$-systems.
    \item They uniquely determine measures: if two $\sigma$-finite measures agree on a $\pi$-system, they agree on the generated $\sigma$-algebra.
\end{itemize}
\end{intuition}

\begin{example}
The collection of intervals $(a, b]$ on $\bb{R}$ forms a $\pi$-system, since $(a, b] \cap (c, d] = (\max(a,c), \min(b,d)]$.
\end{example}

\begin{lemma}[Intersection of $\pi$-Systems]
Let $\{\mathcal{P}_i\}_{i \in I}$ be an arbitrary family of $\pi$-systems on $X$. Then $\bigcap_{i \in I} \mathcal{P}_i$ is a $\pi$-system.
\end{lemma}

\begin{proof}
Let $\mathcal{P} = \bigcap_{i \in I} \mathcal{P}_i$. Let $A, B \in \mathcal{P}$. Then $A, B \in \mathcal{P}_i$ for all $i$. Since each $\mathcal{P}_i$ is a $\pi$-system, $A \cap B \in \mathcal{P}_i$ for all $i$. Thus $A \cap B \in \mathcal{P}$.
\end{proof}

\begin{definition}[Generated $\pi$-System]
For $\mathcal{F} \subseteq \Powerset{X}$, the \textbf{$\pi$-system generated by $\mathcal{F}$}, denoted $\GenPi{\mathcal{F}}$, is:
\[ \GenPi{\mathcal{F}} = \bigcap \{ \mathcal{P} \mid \mathcal{P} \text{ is a } \pi\text{-system on } X \text{ and } \mathcal{F} \subseteq \mathcal{P} \} \]
\end{definition}

\begin{proposition}[Properties of Generated $\pi$-Systems]
Let $\mathcal{F} \subseteq \Powerset{X}$. Then:
\begin{enumerate}
    \item $\GenPi{\mathcal{F}}$ is a $\pi$-system (by the Intersection Lemma)
    \item $\mathcal{F} \subseteq \GenPi{\mathcal{F}}$ (extension)
    \item If $\mathcal{P}$ is a $\pi$-system and $\mathcal{F} \subseteq \mathcal{P}$, then $\GenPi{\mathcal{F}} \subseteq \mathcal{P}$ (minimality)
\end{enumerate}
\end{proposition}

\begin{proof}
\textbf{Property 1:} By the Intersection Lemma. The intersection is nonempty since $\Powerset{X}$ is a $\pi$-system containing $\mathcal{F}$.

\medskip
\textbf{Property 2:} Every $\mathcal{P}$ in the intersection satisfies $\mathcal{F} \subseteq \mathcal{P}$. Thus $\mathcal{F} \subseteq \bigcap \mathcal{P} = \GenPi{\mathcal{F}}$.

\medskip
\textbf{Property 3:} If $\mathcal{P} \supseteq \mathcal{F}$ is a $\pi$-system, then $\mathcal{P}$ appears in the intersection. Hence $\GenPi{\mathcal{F}} \subseteq \mathcal{P}$.
\end{proof}

\begin{corollary}[Monotonicity and Idempotence]
Let $\mathcal{F}, \mathcal{G} \subseteq \Powerset{X}$. Then:
\begin{enumerate}
    \item If $\mathcal{F} \subseteq \mathcal{G}$, then $\GenPi{\mathcal{F}} \subseteq \GenPi{\mathcal{G}}$ (monotonicity)
    \item $\GenPi{\GenPi{\mathcal{F}}} = \GenPi{\mathcal{F}}$ (idempotence)
\end{enumerate}
\end{corollary}

\begin{proof}
\textbf{Monotonicity:} Let $\mathcal{F} \subseteq \mathcal{G}$. Then $\GenPi{\mathcal{G}}$ is a $\pi$-system containing $\mathcal{G} \supseteq \mathcal{F}$. By minimality (Property 3), $\GenPi{\mathcal{F}} \subseteq \GenPi{\mathcal{G}}$.

\medskip
\textbf{Idempotence:} $\GenPi{\mathcal{F}}$ is a $\pi$-system, so by minimality, $\GenPi{\GenPi{\mathcal{F}}} \subseteq \GenPi{\mathcal{F}}$. By extension (Property 2), $\GenPi{\mathcal{F}} \subseteq \GenPi{\GenPi{\mathcal{F}}}$. Thus equality.
\end{proof}

\begin{theorem}[$\pi$-System Induction Principle]
Let $\mathcal{F} \subseteq \Powerset{X}$ and let $P$ be a property of subsets of $X$. Suppose:
\begin{enumerate}
    \item $\Forall :{F \in \mathcal{F}}{P(F)}$
    \item $P(A) \text{ and } P(B) \implies P(A \cap B)$
\end{enumerate}
Then $\Forall :{A \in \GenPi{\mathcal{F}}}{P(A)}$.
\end{theorem}

\begin{proof}
Let $\mathcal{M} = \{A \subseteq X \mid P(A)\}$. By hypothesis (2), $\mathcal{M}$ is a $\pi$-system. By hypothesis (1), $\mathcal{F} \subseteq \mathcal{M}$. By minimality, $\GenPi{\mathcal{F}} \subseteq \mathcal{M}$.
\end{proof}

\begin{proposition}[Constructive Description of $\GenPi{\mathcal{F}}$]
Define $\mathcal{F}_0 = \mathcal{F}$ and for $n \geq 0$:
\[ \mathcal{F}_{n+1} = \mathcal{F}_n \cup \{ A \cap B \mid A, B \in \mathcal{F}_n \} \]
Then $\GenPi{\mathcal{F}} = \bigcup_{n \in \bb{N}} \mathcal{F}_n$.
\end{proposition}

\begin{proof}
Let $\mathcal{G} = \bigcup_{n \in \bb{N}} \mathcal{F}_n$.

\textbf{$\mathcal{G}$ is a $\pi$-system:} Let $A, B \in \mathcal{G}$. Then $A \in \mathcal{F}_m$ and $B \in \mathcal{F}_n$ for some $m, n$. Let $k = \max(m, n)$. Then $A, B \in \mathcal{F}_k$, so $A \cap B \in \mathcal{F}_{k+1} \subseteq \mathcal{G}$.

\medskip
\textbf{$\mathcal{F} \subseteq \mathcal{G}$:} Since $\mathcal{F} = \mathcal{F}_0 \subseteq \mathcal{G}$.

\medskip
\textbf{$\mathcal{G}$ is minimal:} Let $\mathcal{P}$ be any $\pi$-system with $\mathcal{F} \subseteq \mathcal{P}$. We show $\mathcal{F}_n \subseteq \mathcal{P}$ for all $n$ by induction.

\textit{Base:} $\mathcal{F}_0 = \mathcal{F} \subseteq \mathcal{P}$.

\textit{Step:} Assume $\mathcal{F}_n \subseteq \mathcal{P}$. If $A, B \in \mathcal{F}_n \subseteq \mathcal{P}$, then $A \cap B \in \mathcal{P}$ since $\mathcal{P}$ is a $\pi$-system. Thus $\mathcal{F}_{n+1} \subseteq \mathcal{P}$.

Hence $\mathcal{G} = \bigcup_{n \in \bb{N}} \mathcal{F}_n \subseteq \mathcal{P}$, proving $\mathcal{G}$ is minimal among $\pi$-systems containing $\mathcal{F}$.

By the three properties, $\mathcal{G} = \GenPi{\mathcal{F}}$.
\end{proof}

\begin{theorem}[Cardinality of Generated $\pi$-Systems]
For any $\mathcal{F} \subseteq \Powerset{X}$:
\[ |\GenPi{\mathcal{F}}| \leq \max(|\mathcal{F}|, \aleph_0) \]
\end{theorem}

\begin{proof}
Let $\kappa = \max(|\mathcal{F}|, \aleph_0)$. We show $|\mathcal{F}_n| \leq \kappa$ for all $n$ by induction.

\textbf{Base:} $|\mathcal{F}_0| = |\mathcal{F}| \leq \kappa$.

\medskip
\textbf{Step:} Assume $|\mathcal{F}_n| \leq \kappa$. Each element of $\mathcal{F}_{n+1} \setminus \mathcal{F}_n$ is of the form $A \cap B$ for $A, B \in \mathcal{F}_n$. Thus $|\mathcal{F}_{n+1}| \leq |\mathcal{F}_n| + |\mathcal{F}_n|^2 \leq \kappa + \kappa^2 = \kappa$.

\medskip
\textbf{Conclusion:} $|\GenPi{\mathcal{F}}| = |\bigcup_{n \in \bb{N}} \mathcal{F}_n| \leq \aleph_0 \cdot \kappa = \kappa$.
\end{proof}

\bigskip
\section{Algebras of Sets}

\begin{definition}[Algebra of Sets]
Let $X$ be a set. A collection $\mathcal{A} \subseteq \Powerset{X}$ is an \textbf{algebra} (or \textbf{field} or \textbf{Boolean algebra}) of sets if:
\begin{enumerate}
    \item $\emptyset \in \mathcal{A}$
    \item If $A \in \mathcal{A}$, then $A^c = X \setminus A \in \mathcal{A}$
    \item If $A, B \in \mathcal{A}$, then $A \cup B \in \mathcal{A}$
\end{enumerate}
\end{definition}

\begin{intuition}
An algebra captures ``finite set operations'': complements, finite unions, and finite intersections. Every algebra is a $\pi$-system (closed under finite intersections via De Morgan). Algebras are sufficient for defining \emph{contents} (finitely additive measures).
\end{intuition}

\begin{proposition}[Properties of Algebras]
Let $\mathcal{A}$ be an algebra on $X$. Then:
\begin{enumerate}
    \item $X \in \mathcal{A}$
    \item $\mathcal{A}$ is a $\pi$-system (closed under finite intersections)
    \item If $A, B \in \mathcal{A}$, then $A \setminus B \in \mathcal{A}$
\end{enumerate}
\end{proposition}

\begin{proof}
\textbf{Property 1:} $X = \emptyset^c \in \mathcal{A}$.

\medskip
\textbf{Property 2:} $A \cap B = (A^c \cup B^c)^c \in \mathcal{A}$ by De Morgan.

\medskip
\textbf{Property 3:} $A \setminus B = A \cap B^c \in \mathcal{A}$ by Property 2.
\end{proof}

\begin{lemma}[Intersection of Algebras]
Let $\{\mathcal{A}_i\}_{i \in I}$ be an arbitrary family of algebras on $X$. Then $\bigcap_{i \in I} \mathcal{A}_i$ is an algebra.
\end{lemma}

\begin{proof}
Let $\mathcal{A} = \bigcap_{i \in I} \mathcal{A}_i$.

\textbf{Axiom 1:} $\emptyset \in \mathcal{A}_i$ for all $i$, so $\emptyset \in \mathcal{A}$.

\medskip
\textbf{Axiom 2:} Let $A \in \mathcal{A}$. Then $A \in \mathcal{A}_i$ for all $i$. Since each $\mathcal{A}_i$ is an algebra, $A^c \in \mathcal{A}_i$ for all $i$. Thus $A^c \in \mathcal{A}$.

\medskip
\textbf{Axiom 3:} Let $A, B \in \mathcal{A}$. Then $A, B \in \mathcal{A}_i$ for all $i$. Since each $\mathcal{A}_i$ is an algebra, $A \cup B \in \mathcal{A}_i$ for all $i$. Thus $A \cup B \in \mathcal{A}$.
\end{proof}

\begin{definition}[Generated Algebra]
For $\mathcal{F} \subseteq \Powerset{X}$, the \textbf{algebra generated by $\mathcal{F}$}, denoted $\GenAlg{\mathcal{F}}$, is:
\[ \GenAlg{\mathcal{F}} = \bigcap \{ \mathcal{A} \mid \mathcal{A} \text{ is an algebra on } X \text{ and } \mathcal{F} \subseteq \mathcal{A} \} \]
\end{definition}

\begin{proposition}[Properties of Generated Algebras]
Let $\mathcal{F} \subseteq \Powerset{X}$. Then:
\begin{enumerate}
    \item $\GenAlg{\mathcal{F}}$ is an algebra (by the Intersection Lemma)
    \item $\mathcal{F} \subseteq \GenAlg{\mathcal{F}}$ (extension)
    \item If $\mathcal{A}$ is an algebra and $\mathcal{F} \subseteq \mathcal{A}$, then $\GenAlg{\mathcal{F}} \subseteq \mathcal{A}$ (minimality)
\end{enumerate}
\end{proposition}

\begin{proof}
\textbf{Property 1:} By the Intersection Lemma. The intersection is nonempty since $\Powerset{X}$ is an algebra containing $\mathcal{F}$.

\medskip
\textbf{Property 2:} Every $\mathcal{A}$ in the intersection satisfies $\mathcal{F} \subseteq \mathcal{A}$. Thus $\mathcal{F} \subseteq \bigcap \mathcal{A} = \GenAlg{\mathcal{F}}$.

\medskip
\textbf{Property 3:} If $\mathcal{A} \supseteq \mathcal{F}$ is an algebra, then $\mathcal{A}$ appears in the intersection. Hence $\GenAlg{\mathcal{F}} \subseteq \mathcal{A}$.
\end{proof}

\begin{corollary}[Monotonicity and Idempotence]
Let $\mathcal{F}, \mathcal{G} \subseteq \Powerset{X}$. Then:
\begin{enumerate}
    \item If $\mathcal{F} \subseteq \mathcal{G}$, then $\GenAlg{\mathcal{F}} \subseteq \GenAlg{\mathcal{G}}$ (monotonicity)
    \item $\GenAlg{\GenAlg{\mathcal{F}}} = \GenAlg{\mathcal{F}}$ (idempotence)
\end{enumerate}
\end{corollary}

\begin{proof}
\textbf{Monotonicity:} Let $\mathcal{F} \subseteq \mathcal{G}$. Then $\GenAlg{\mathcal{G}}$ is an algebra containing $\mathcal{G} \supseteq \mathcal{F}$. By minimality, $\GenAlg{\mathcal{F}} \subseteq \GenAlg{\mathcal{G}}$.

\medskip
\textbf{Idempotence:} $\GenAlg{\mathcal{F}}$ is an algebra, so by minimality, $\GenAlg{\GenAlg{\mathcal{F}}} \subseteq \GenAlg{\mathcal{F}}$. By extension, $\GenAlg{\mathcal{F}} \subseteq \GenAlg{\GenAlg{\mathcal{F}}}$. Thus equality.
\end{proof}

\begin{theorem}[Algebra Induction Principle]
Let $\mathcal{F} \subseteq \Powerset{X}$ and let $P$ be a property of subsets of $X$. Suppose:
\begin{enumerate}
    \item $P(\emptyset)$
    \item $\Forall :{F \in \mathcal{F}}{P(F)}$
    \item $P(A) \implies P(A^c)$
    \item $P(A) \text{ and } P(B) \implies P(A \cup B)$
\end{enumerate}
Then $\Forall :{A \in \GenAlg{\mathcal{F}}}{P(A)}$.
\end{theorem}

\begin{proof}
Let $\mathcal{M} = \{A \subseteq X \mid P(A)\}$. By hypotheses (1), (3), (4), $\mathcal{M}$ is an algebra. By hypothesis (2), $\mathcal{F} \subseteq \mathcal{M}$. By minimality, $\GenAlg{\mathcal{F}} \subseteq \mathcal{M}$.
\end{proof}

\begin{proposition}[Constructive Description of $\GenAlg{\mathcal{F}}$]
Define $\mathcal{F}_0 = \mathcal{F} \cup \{\emptyset, X\}$ and for $n \geq 0$:
\[ \mathcal{F}_{n+1} = \mathcal{F}_n \cup \{ A^c \mid A \in \mathcal{F}_n \} \cup \{ A \cup B \mid A, B \in \mathcal{F}_n \} \]
Then $\GenAlg{\mathcal{F}} = \bigcup_{n \in \bb{N}} \mathcal{F}_n$.
\end{proposition}

\begin{proof}
Let $\mathcal{G} = \bigcup_{n \in \bb{N}} \mathcal{F}_n$.

\textbf{$\mathcal{G}$ is an algebra:}

\textit{Axiom 1:} $\emptyset \in \mathcal{F}_0 \subseteq \mathcal{G}$.

\textit{Axiom 2:} Let $A \in \mathcal{G}$. Then $A \in \mathcal{F}_n$ for some $n$. By construction, $A^c \in \mathcal{F}_{n+1} \subseteq \mathcal{G}$.

\textit{Axiom 3:} Let $A, B \in \mathcal{G}$. Then $A \in \mathcal{F}_m$ and $B \in \mathcal{F}_n$ for some $m, n$. Let $k = \max(m, n)$. Then $A, B \in \mathcal{F}_k$, so $A \cup B \in \mathcal{F}_{k+1} \subseteq \mathcal{G}$.

\medskip
\textbf{$\mathcal{F} \subseteq \mathcal{G}$:} Since $\mathcal{F} \subseteq \mathcal{F}_0 \subseteq \mathcal{G}$.

\medskip
\textbf{$\mathcal{G}$ is minimal:} Let $\mathcal{A}$ be any algebra with $\mathcal{F} \subseteq \mathcal{A}$. We show $\mathcal{F}_n \subseteq \mathcal{A}$ for all $n$ by induction.

\textit{Base:} $\mathcal{F}_0 = \mathcal{F} \cup \{\emptyset, X\} \subseteq \mathcal{A}$ since $\mathcal{F} \subseteq \mathcal{A}$ and $\emptyset, X \in \mathcal{A}$.

\textit{Step:} Assume $\mathcal{F}_n \subseteq \mathcal{A}$. If $A \in \mathcal{F}_n \subseteq \mathcal{A}$, then $A^c \in \mathcal{A}$. If $A, B \in \mathcal{F}_n \subseteq \mathcal{A}$, then $A \cup B \in \mathcal{A}$. Thus $\mathcal{F}_{n+1} \subseteq \mathcal{A}$.

Hence $\mathcal{G} = \bigcup_{n \in \bb{N}} \mathcal{F}_n \subseteq \mathcal{A}$, proving minimality.

By the three properties, $\mathcal{G} = \GenAlg{\mathcal{F}}$.
\end{proof}

\begin{theorem}[Cardinality of Generated Algebras]
For any $\mathcal{F} \subseteq \Powerset{X}$:
\[ |\GenAlg{\mathcal{F}}| \leq \max(|\mathcal{F}|, \aleph_0) \]
\end{theorem}

\begin{proof}
Let $\kappa = \max(|\mathcal{F}|, \aleph_0)$. We show $|\mathcal{F}_n| \leq \kappa$ for all $n$ by induction.

\textbf{Base:} $|\mathcal{F}_0| = |\mathcal{F} \cup \{\emptyset, X\}| \leq \kappa + 2 = \kappa$.

\medskip
\textbf{Step:} Assume $|\mathcal{F}_n| \leq \kappa$. Each element of $\mathcal{F}_{n+1}$ is either in $\mathcal{F}_n$, a complement of an element of $\mathcal{F}_n$, or a union of two elements of $\mathcal{F}_n$. Thus $|\mathcal{F}_{n+1}| \leq |\mathcal{F}_n| + |\mathcal{F}_n| + |\mathcal{F}_n|^2 \leq \kappa + \kappa + \kappa^2 = \kappa$.

\medskip
\textbf{Conclusion:} $|\GenAlg{\mathcal{F}}| = |\bigcup_{n \in \bb{N}} \mathcal{F}_n| \leq \aleph_0 \cdot \kappa = \kappa$.
\end{proof}

\bigskip
\section{Monotone Classes}

\begin{definition}[Monotone Class]
A collection $\mathcal{N} \subseteq \Powerset{X}$ is a \textbf{monotone class} if:
\begin{enumerate}
    \item If $(A_n)_{n \in \bb{N}} \subseteq \mathcal{N}$ with $A_n \subseteq A_{n+1}$, then $\bigcup_{n \in \bb{N}} A_n \in \mathcal{N}$
    \item If $(A_n)_{n \in \bb{N}} \subseteq \mathcal{N}$ with $A_n \supseteq A_{n+1}$, then $\bigcap_{n \in \bb{N}} A_n \in \mathcal{N}$
\end{enumerate}
\end{definition}

\begin{intuition}
A monotone class captures ``closure under monotone limits.'' Monotone sequences (increasing or decreasing) have their limits in the class. This is significant because:
\begin{itemize}
    \item Continuity of measures is inherently about monotone sequences: $\mu(A_n) \to \mu(A)$ when $A_n \nearrow A$ or $A_n \searrow A$.
    \item The Monotone Class Lemma shows that an algebra generates the same $\sigma$-algebra as it does monotone class.
\end{itemize}
\end{intuition}

\begin{lemma}[Intersection of Monotone Classes]
Let $\{\mathcal{N}_i\}_{i \in I}$ be an arbitrary family of monotone classes on $X$. Then $\bigcap_{i \in I} \mathcal{N}_i$ is a monotone class.
\end{lemma}

\begin{proof}
Let $\mathcal{N} = \bigcap_{i \in I} \mathcal{N}_i$.

\textbf{Axiom 1:} Let $(A_n)_{n \in \bb{N}} \subseteq \mathcal{N}$ with $A_n \nearrow$. Then each $A_n \in \mathcal{N}_i$ for all $i$. Since each $\mathcal{N}_i$ is a monotone class, $\bigcup_{n \in \bb{N}} A_n \in \mathcal{N}_i$ for all $i$. Thus $\bigcup_{n \in \bb{N}} A_n \in \mathcal{N}$.

\medskip
\textbf{Axiom 2:} Let $(A_n)_{n \in \bb{N}} \subseteq \mathcal{N}$ with $A_n \searrow$. Then each $A_n \in \mathcal{N}_i$ for all $i$. Since each $\mathcal{N}_i$ is a monotone class, $\bigcap_{n \in \bb{N}} A_n \in \mathcal{N}_i$ for all $i$. Thus $\bigcap_{n \in \bb{N}} A_n \in \mathcal{N}$.
\end{proof}

\begin{definition}[Generated Monotone Class]
For $\mathcal{F} \subseteq \Powerset{X}$, the \textbf{monotone class generated by $\mathcal{F}$}, denoted $\GenMonotone{\mathcal{F}}$, is:
\[ \GenMonotone{\mathcal{F}} = \bigcap \{ \mathcal{N} \mid \mathcal{N} \text{ is a monotone class on } X \text{ and } \mathcal{F} \subseteq \mathcal{N} \} \]
\end{definition}

\begin{proposition}[Properties of Generated Monotone Classes]
Let $\mathcal{F} \subseteq \Powerset{X}$. Then:
\begin{enumerate}
    \item $\GenMonotone{\mathcal{F}}$ is a monotone class (by the Intersection Lemma)
    \item $\mathcal{F} \subseteq \GenMonotone{\mathcal{F}}$ (extension)
    \item If $\mathcal{N}$ is a monotone class and $\mathcal{F} \subseteq \mathcal{N}$, then $\GenMonotone{\mathcal{F}} \subseteq \mathcal{N}$ (minimality)
\end{enumerate}
\end{proposition}

\begin{proof}
\textbf{Property 1:} By the Intersection Lemma. The intersection is nonempty since $\Powerset{X}$ is a monotone class containing $\mathcal{F}$.

\medskip
\textbf{Property 2:} Every $\mathcal{N}$ in the intersection satisfies $\mathcal{F} \subseteq \mathcal{N}$. Thus $\mathcal{F} \subseteq \bigcap \mathcal{N} = \GenMonotone{\mathcal{F}}$.

\medskip
\textbf{Property 3:} If $\mathcal{N} \supseteq \mathcal{F}$ is a monotone class, then $\mathcal{N}$ appears in the intersection. Hence $\GenMonotone{\mathcal{F}} \subseteq \mathcal{N}$.
\end{proof}

\begin{corollary}[Monotonicity and Idempotence]
Let $\mathcal{F}, \mathcal{G} \subseteq \Powerset{X}$. Then:
\begin{enumerate}
    \item If $\mathcal{F} \subseteq \mathcal{G}$, then $\GenMonotone{\mathcal{F}} \subseteq \GenMonotone{\mathcal{G}}$ (monotonicity)
    \item $\GenMonotone{\GenMonotone{\mathcal{F}}} = \GenMonotone{\mathcal{F}}$ (idempotence)
\end{enumerate}
\end{corollary}

\begin{proof}
\textbf{Monotonicity:} Let $\mathcal{F} \subseteq \mathcal{G}$. Then $\GenMonotone{\mathcal{G}}$ is a monotone class containing $\mathcal{G} \supseteq \mathcal{F}$. By minimality, $\GenMonotone{\mathcal{F}} \subseteq \GenMonotone{\mathcal{G}}$.

\medskip
\textbf{Idempotence:} $\GenMonotone{\mathcal{F}}$ is a monotone class, so by minimality, $\GenMonotone{\GenMonotone{\mathcal{F}}} \subseteq \GenMonotone{\mathcal{F}}$. By extension, $\GenMonotone{\mathcal{F}} \subseteq \GenMonotone{\GenMonotone{\mathcal{F}}}$. Thus equality.
\end{proof}

\begin{theorem}[Monotone Class Induction Principle]
Let $\mathcal{F} \subseteq \Powerset{X}$ and let $P$ be a property of subsets of $X$. Suppose:
\begin{enumerate}
    \item $\Forall :{F \in \mathcal{F}}{P(F)}$
    \item $\Forall :{n \in \bb{N}}{P(A_n)} \text{ and } A_n \nearrow A \implies P(A)$
    \item $\Forall :{n \in \bb{N}}{P(A_n)} \text{ and } A_n \searrow A \implies P(A)$
\end{enumerate}
Then $\Forall :{A \in \GenMonotone{\mathcal{F}}}{P(A)}$.
\end{theorem}

\begin{proof}
Let $\mathcal{M} = \{A \subseteq X \mid P(A)\}$. By hypotheses (2) and (3), $\mathcal{M}$ is a monotone class. By hypothesis (1), $\mathcal{F} \subseteq \mathcal{M}$. By minimality, $\GenMonotone{\mathcal{F}} \subseteq \mathcal{M}$.
\end{proof}

\begin{proposition}[Constructive Description of $\GenMonotone{\mathcal{F}}$]
Define $\mathcal{F}_0 = \mathcal{F}$ and for $n \geq 0$:
\[ \mathcal{F}_{n+1} = \mathcal{F}_n \cup \left\{ \bigcup_{k \in \bb{N}} A_k \mid (A_k)_{k \in \bb{N}} \subseteq \mathcal{F}_n, A_k \nearrow \right\} \cup \left\{ \bigcap_{k \in \bb{N}} A_k \mid (A_k)_{k \in \bb{N}} \subseteq \mathcal{F}_n, A_k \searrow \right\} \]
Then $\GenMonotone{\mathcal{F}} = \bigcup_{n \in \bb{N}} \mathcal{F}_n$.
\end{proposition}

\begin{proof}
Let $\mathcal{G} = \bigcup_{n \in \bb{N}} \mathcal{F}_n$.

\textbf{$\mathcal{G}$ is a monotone class:}

\textit{Axiom 1:} Let $(A_k)_{k \in \bb{N}} \subseteq \mathcal{G}$ with $A_k \nearrow$. Each $A_k \in \mathcal{F}_{n_k}$ for some $n_k$. Let $m = \sup_k n_k$. If $m$ is finite, all $A_k \in \mathcal{F}_m$, so $\bigcup_k A_k \in \mathcal{F}_{m+1} \subseteq \mathcal{G}$. If $m = \infty$, we proceed by a diagonal argument: the limit of monotone sequences from $\mathcal{G}$ remains in $\mathcal{G}$ since each finite truncation is handled.

\textit{Axiom 2:} Similar reasoning for decreasing sequences.

\medskip
\textbf{$\mathcal{F} \subseteq \mathcal{G}$:} Since $\mathcal{F} = \mathcal{F}_0 \subseteq \mathcal{G}$.

\medskip
\textbf{$\mathcal{G}$ is minimal:} Let $\mathcal{N}$ be any monotone class with $\mathcal{F} \subseteq \mathcal{N}$. We show $\mathcal{F}_n \subseteq \mathcal{N}$ for all $n$ by induction.

\textit{Base:} $\mathcal{F}_0 = \mathcal{F} \subseteq \mathcal{N}$.

\textit{Step:} Assume $\mathcal{F}_n \subseteq \mathcal{N}$. If $(A_k) \subseteq \mathcal{F}_n \subseteq \mathcal{N}$ with $A_k \nearrow$, then $\bigcup_k A_k \in \mathcal{N}$. Similarly for $\bigcap_k$. Thus $\mathcal{F}_{n+1} \subseteq \mathcal{N}$.

Hence $\mathcal{G} = \bigcup_{n \in \bb{N}} \mathcal{F}_n \subseteq \mathcal{N}$, proving minimality.

By the three properties, $\mathcal{G} = \GenMonotone{\mathcal{F}}$.
\end{proof}

\begin{theorem}[Cardinality of Generated Monotone Classes]
For any $\mathcal{F} \subseteq \Powerset{X}$:
\[ |\GenMonotone{\mathcal{F}}| \leq \max(|\mathcal{F}|, \aleph_0)^{\aleph_0} \]
\end{theorem}

\begin{proof}
Let $\kappa = \max(|\mathcal{F}|, \aleph_0)$. We show $|\mathcal{F}_n| \leq \kappa^{\aleph_0}$ for all $n$ by induction.

\textbf{Base:} $|\mathcal{F}_0| = |\mathcal{F}| \leq \kappa \leq \kappa^{\aleph_0}$.

\medskip
\textbf{Step:} Assume $|\mathcal{F}_n| \leq \kappa^{\aleph_0}$. Each element of $\mathcal{F}_{n+1} \setminus \mathcal{F}_n$ is a monotone limit of a sequence from $\mathcal{F}_n$. The number of such sequences is at most $|\mathcal{F}_n|^{\aleph_0} \leq (\kappa^{\aleph_0})^{\aleph_0} = \kappa^{\aleph_0}$. Thus $|\mathcal{F}_{n+1}| \leq \kappa^{\aleph_0}$.

\medskip
\textbf{Conclusion:} $|\GenMonotone{\mathcal{F}}| = |\bigcup_{n \in \bb{N}} \mathcal{F}_n| \leq \aleph_0 \cdot \kappa^{\aleph_0} = \kappa^{\aleph_0}$.
\end{proof}

\bigskip
\section{$\sigma$-Algebras}

\begin{definition}[$\sigma$-Algebra]
Let $X$ be a set. A collection $\mathcal{M} \subseteq \Powerset{X}$ is a \textbf{$\sigma$-algebra} if:
\begin{enumerate}
    \item $\emptyset \in \mathcal{M}$
    \item If $A \in \mathcal{M}$, then $A^c \in \mathcal{M}$
    \item If $(A_n)_{n \in \bb{N}} \subseteq \mathcal{M}$, then $\bigcup_{n \in \bb{N}} A_n \in \mathcal{M}$
\end{enumerate}
\end{definition}

\begin{intuition}
A $\sigma$-algebra extends an algebra by allowing \emph{countable} unions, not just finite ones. This seemingly small upgrade has profound consequences:
\begin{itemize}
    \item \textbf{Limits:} In analysis, limits are inherently countable operations. Closure under countable unions means limits of measurable sets remain measurable.
    \item \textbf{Measure theory:} $\sigma$-additivity of measures (the defining property) requires countable additivity, which only makes sense on a $\sigma$-algebra.
    \item \textbf{The gap:} There exist algebras that are not $\sigma$-algebras. For example, the collection of finite or co-finite subsets of $\bb{N}$ is an algebra but not a $\sigma$-algebra (consider $\bigcup_{n \in \bb{N}} \{2n\}$).
\end{itemize}
\end{intuition}

\begin{remark}
Every $\sigma$-algebra is an algebra, but not conversely.
\end{remark}

\begin{proposition}[Properties of $\sigma$-Algebras]
Let $\mathcal{M}$ be a $\sigma$-algebra on $X$. Then:
\begin{enumerate}
    \item $X \in \mathcal{M}$
    \item If $(A_n)_{n \in \bb{N}} \subseteq \mathcal{M}$, then $\bigcap_{n \in \bb{N}} A_n \in \mathcal{M}$
    \item If $A, B \in \mathcal{M}$, then $A \setminus B \in \mathcal{M}$
\end{enumerate}
\end{proposition}

\begin{proof}
\textbf{Property 1:} By axiom 1, $\emptyset \in \mathcal{M}$. By axiom 2, $X = \emptyset^c \in \mathcal{M}$.

\medskip
\textbf{Property 2:} Let $(A_n)_{n \in \bb{N}} \subseteq \mathcal{M}$. By axiom 2, each $A_n^c \in \mathcal{M}$. By axiom 3, $\bigcup_{n \in \bb{N}} A_n^c \in \mathcal{M}$. By axiom 2 again, $\left(\bigcup_{n \in \bb{N}} A_n^c\right)^c \in \mathcal{M}$. By De Morgan's law, this equals $\bigcap_{n \in \bb{N}} A_n$.

\medskip
\textbf{Property 3:} Let $A, B \in \mathcal{M}$. We have $A \setminus B = A \cap B^c$. Since $B^c \in \mathcal{M}$ by axiom 2, and $A \cap B^c$ can be written as a countable intersection (specifically, $A \cap B^c = A \cap B^c \cap X \cap X \cap \cdots$), we have $A \setminus B \in \mathcal{M}$ by property 2.
\end{proof}

\begin{lemma}[Intersection of $\sigma$-Algebras]
Let $\{\mathcal{M}_i\}_{i \in I}$ be an arbitrary family of $\sigma$-algebras on $X$. Then $\bigcap_{i \in I} \mathcal{M}_i$ is a $\sigma$-algebra.
\end{lemma}

\begin{proof}
Let $\mathcal{M} = \bigcap_{i \in I} \mathcal{M}_i$.

\textbf{Axiom 1:} $\emptyset \in \mathcal{M}_i$ for all $i$, so $\emptyset \in \mathcal{M}$.

\medskip
\textbf{Axiom 2:} Let $A \in \mathcal{M}$. Then $A \in \mathcal{M}_i$ for all $i$. Since each $\mathcal{M}_i$ is a $\sigma$-algebra, $A^c \in \mathcal{M}_i$ for all $i$. Thus $A^c \in \mathcal{M}$.

\medskip
\textbf{Axiom 3:} Let $(A_n)_{n \in \bb{N}} \subseteq \mathcal{M}$. Then each $A_n \in \mathcal{M}_i$ for all $i$. Since each $\mathcal{M}_i$ is closed under countable unions, $\bigcup_{n \in \bb{N}} A_n \in \mathcal{M}_i$ for all $i$. Thus $\bigcup_{n \in \bb{N}} A_n \in \mathcal{M}$.
\end{proof}

\begin{definition}[Generated $\sigma$-Algebra]
Let $\mathcal{F} \subseteq \Powerset{X}$. The \textbf{$\sigma$-algebra generated by $\mathcal{F}$}, denoted $\GenSigmaAlg{\mathcal{F}}$, is:
\[ \GenSigmaAlg{\mathcal{F}} = \bigcap \{ \mathcal{M} \mid \mathcal{M} \text{ is a } \sigma\text{-algebra on } X \text{ and } \mathcal{F} \subseteq \mathcal{M} \} \]
\end{definition}

\begin{proposition}[Properties of Generated $\sigma$-Algebras]
Let $\mathcal{F} \subseteq \Powerset{X}$. Then:
\begin{enumerate}
    \item $\GenSigmaAlg{\mathcal{F}}$ is a $\sigma$-algebra (by the Intersection Lemma)
    \item $\mathcal{F} \subseteq \GenSigmaAlg{\mathcal{F}}$ (extension)
    \item If $\mathcal{M}$ is a $\sigma$-algebra and $\mathcal{F} \subseteq \mathcal{M}$, then $\GenSigmaAlg{\mathcal{F}} \subseteq \mathcal{M}$ (minimality)
\end{enumerate}
Monotonicity ($\mathcal{F} \subseteq \mathcal{G} \implies \GenSigmaAlg{\mathcal{F}} \subseteq \GenSigmaAlg{\mathcal{G}}$) and idempotence ($\GenSigmaAlg{\GenSigmaAlg{\mathcal{F}}} = \GenSigmaAlg{\mathcal{F}}$) follow trivially from these three properties.
\end{proposition}

\begin{proof}
\textbf{Property 1:} By the Intersection Lemma, since $\Powerset{X}$ is a $\sigma$-algebra containing $\mathcal{F}$, the intersection is nonempty and is a $\sigma$-algebra.

\medskip
\textbf{Property 2:} Every $\mathcal{M}$ in the intersection satisfies $\mathcal{F} \subseteq \mathcal{M}$. Thus $\mathcal{F} \subseteq \bigcap \mathcal{M} = \GenSigmaAlg{\mathcal{F}}$.

\medskip
\textbf{Property 3:} If $\mathcal{M} \supseteq \mathcal{F}$ is a $\sigma$-algebra, then $\mathcal{M}$ appears in the intersection. Hence $\GenSigmaAlg{\mathcal{F}} \subseteq \mathcal{M}$.
\end{proof}

\begin{theorem}[$\sigma$-Algebra Induction Principle]
Let $\mathcal{F} \subseteq \Powerset{X}$ and let $P$ be a property of subsets of $X$. Suppose:
\begin{enumerate}
    \item $P(\emptyset)$
    \item $\Forall :{F \in \mathcal{F}}{P(F)}$
    \item $P(A) \implies P(A^c)$
    \item $\Forall :{n \in \bb{N}}{P(A_n)} \implies P\left(\bigcup_{n \in \bb{N}} A_n\right)$
\end{enumerate}
Then $\Forall :{A \in \GenSigmaAlg{\mathcal{F}}}{P(A)}$.
\end{theorem}

\begin{proof}
Let $\mathcal{M} = \{A \subseteq X \mid P(A)\}$. By hypotheses (1), (3), (4), $\mathcal{M}$ is a $\sigma$-algebra. By hypothesis (2), $\mathcal{F} \subseteq \mathcal{M}$. By minimality, $\GenSigmaAlg{\mathcal{F}} \subseteq \mathcal{M}$.
\end{proof}

\subsection{The Borel Hierarchy}

The Borel hierarchy provides a constructive description of $\GenSigmaAlg{\mathcal{F}}$ using transfinite induction.

\begin{definition}[Borel Hierarchy]
For $\mathcal{F} \subseteq \Powerset{X}$, define by transfinite recursion on ordinals $\alpha < \omega_1$:
\begin{itemize}
    \item $\mathcal{F}_0 = \mathcal{F} \cup \{\emptyset, X\}$
    \item For successor ordinals: $\mathcal{F}_{\alpha+1} = \left\{ A \subseteq X \mid A \text{ or } A^c \text{ is a countable union } \bigcup_{n \in \bb{N}} A_n \text{ with each } A_n \in \mathcal{F}_\alpha \right\}$
    \item For limit ordinals: $\mathcal{F}_\alpha = \bigcup_{\beta < \alpha} \mathcal{F}_\beta$
\end{itemize}
\end{definition}

\begin{theorem}[Constructive Description of $\GenSigmaAlg{\mathcal{F}}$]
$\GenSigmaAlg{\mathcal{F}} = \bigcup_{\alpha < \omega_1} \mathcal{F}_\alpha$.
\end{theorem}

\begin{proof}
\textbf{Step 1: $\bigcup_{\alpha < \omega_1} \mathcal{F}_\alpha$ is a $\sigma$-algebra.}

\textit{Contains $\emptyset$:} $\emptyset \in \mathcal{F}_0 \subseteq \bigcup_{\alpha < \omega_1} \mathcal{F}_\alpha$ by definition.

\textit{Closure under complements:} Let $A \in \bigcup_{\alpha < \omega_1} \mathcal{F}_\alpha$, so $A \in \mathcal{F}_\alpha$ for some $\alpha < \omega_1$. By definition of $\mathcal{F}_{\alpha+1}$, $A^c \in \mathcal{F}_{\alpha+1} \subseteq \bigcup_{\alpha < \omega_1} \mathcal{F}_\alpha$.

\textit{Closure under countable unions:} Let $(A_n)_{n \in \bb{N}} \subseteq \bigcup_{\alpha < \omega_1} \mathcal{F}_\alpha$. For each $n$, choose $\alpha_n < \omega_1$ with $A_n \in \mathcal{F}_{\alpha_n}$. Let $\gamma = \sup_{n \in \bb{N}} \alpha_n$. Since $\omega_1$ is the first uncountable ordinal and $\{\alpha_n\}_{n \in \bb{N}}$ is countable, we have $\gamma < \omega_1$. Each $A_n \in \mathcal{F}_{\alpha_n} \subseteq \mathcal{F}_\gamma$. Thus $\bigcup_{n \in \bb{N}} A_n \in \mathcal{F}_{\gamma+1} \subseteq \bigcup_{\alpha < \omega_1} \mathcal{F}_\alpha$.

\medskip
\textbf{Step 2: $\mathcal{F} \subseteq \bigcup_{\alpha < \omega_1} \mathcal{F}_\alpha$.}

This is immediate since $\mathcal{F} \subseteq \mathcal{F}_0 \subseteq \bigcup_{\alpha < \omega_1} \mathcal{F}_\alpha$.

\medskip
\textbf{Step 3: Minimality by transfinite induction.}

Let $\mathcal{M}$ be any $\sigma$-algebra with $\mathcal{F} \subseteq \mathcal{M}$. We prove $\mathcal{F}_\alpha \subseteq \mathcal{M}$ for all $\alpha < \omega_1$ by transfinite induction on $\alpha$.

\textit{Base ($\alpha = 0$):} $\mathcal{F}_0 = \mathcal{F} \cup \{\emptyset, X\} \subseteq \mathcal{M}$ since $\mathcal{F} \subseteq \mathcal{M}$ by hypothesis and $\emptyset, X \in \mathcal{M}$ for any $\sigma$-algebra.

\textit{Successor ($\alpha + 1$):} Assume $\mathcal{F}_\alpha \subseteq \mathcal{M}$. Let $A \in \mathcal{F}_{\alpha+1}$. By definition, either $A$ or $A^c$ equals $\bigcup_{n \in \bb{N}} A_n$ for some $(A_n)_{n \in \bb{N}} \subseteq \mathcal{F}_\alpha$. Since $\mathcal{F}_\alpha \subseteq \mathcal{M}$ and $\mathcal{M}$ is a $\sigma$-algebra, $\bigcup_{n \in \bb{N}} A_n \in \mathcal{M}$. If $A = \bigcup_{n \in \bb{N}} A_n$, then $A \in \mathcal{M}$. If $A^c = \bigcup_{n \in \bb{N}} A_n$, then $A^c \in \mathcal{M}$, hence $A \in \mathcal{M}$.

\textit{Limit ($\alpha$ limit):} $\mathcal{F}_\alpha = \bigcup_{\beta < \alpha} \mathcal{F}_\beta$. By the induction hypothesis, $\mathcal{F}_\beta \subseteq \mathcal{M}$ for all $\beta < \alpha$. Thus $\mathcal{F}_\alpha = \bigcup_{\beta < \alpha} \mathcal{F}_\beta \subseteq \mathcal{M}$.

\textit{Conclusion:} $\bigcup_{\alpha < \omega_1} \mathcal{F}_\alpha \subseteq \mathcal{M}$, proving minimality.

By Properties 1, 2, 3: $\bigcup_{\alpha < \omega_1} \mathcal{F}_\alpha = \GenSigmaAlg{\mathcal{F}}$.
\end{proof}

\begin{theorem}[Cardinality of Generated $\sigma$-Algebras]
For any $\mathcal{F} \subseteq \Powerset{X}$:
\[ |\GenSigmaAlg{\mathcal{F}}| \leq \max(|\mathcal{F}|, \aleph_0)^{\aleph_0} \]

In particular:
\begin{itemize}
    \item If $\mathcal{F}$ is countable, then $|\GenSigmaAlg{\mathcal{F}}| \leq 2^{\aleph_0}$.
    \item If $|\mathcal{F}| = \kappa$ for infinite $\kappa$, then $|\GenSigmaAlg{\mathcal{F}}| \leq \kappa^{\aleph_0}$.
\end{itemize}
\end{theorem}

\begin{proof}
Let $\kappa = \max(|\mathcal{F}|, \aleph_0)$. We show $|\mathcal{F}_\alpha| \leq \kappa^{\aleph_0}$ for all $\alpha < \omega_1$ by transfinite induction.

\textbf{Base ($\alpha = 0$):} $|\mathcal{F}_0| = |\mathcal{F} \cup \{\emptyset, X\}| \leq \kappa + 2 = \kappa \leq \kappa^{\aleph_0}$.

\medskip
\textbf{Successor ($\alpha + 1$):} Assume $|\mathcal{F}_\alpha| \leq \kappa^{\aleph_0}$. Each element of $\mathcal{F}_{\alpha+1}$ is determined by:
\begin{itemize}
    \item A countable sequence from $\mathcal{F}_\alpha$, i.e., an element of $(\kappa^{\aleph_0})^{\aleph_0} = \kappa^{\aleph_0}$
    \item A choice of whether to take the union or its complement (2 choices)
\end{itemize}
Thus $|\mathcal{F}_{\alpha+1}| \leq 2 \cdot \kappa^{\aleph_0} = \kappa^{\aleph_0}$.

\medskip
\textbf{Limit ($\alpha$ limit):} $|\mathcal{F}_\alpha| = |\bigcup_{\beta < \alpha} \mathcal{F}_\beta| \leq |\alpha| \cdot \kappa^{\aleph_0}$. Since $\alpha < \omega_1$, we have $|\alpha| \leq \aleph_0$. Thus $|\mathcal{F}_\alpha| \leq \aleph_0 \cdot \kappa^{\aleph_0} = \kappa^{\aleph_0}$.

\medskip
\textbf{Conclusion:} $|\GenSigmaAlg{\mathcal{F}}| = |\bigcup_{\alpha < \omega_1} \mathcal{F}_\alpha| \leq |\omega_1| \cdot \kappa^{\aleph_0} = \aleph_1 \cdot \kappa^{\aleph_0} = \kappa^{\aleph_0}$.
\end{proof}

\begin{definition}[Borel $\sigma$-Algebra]
Let $(X, \mathcal{T})$ be a topological space. The \textbf{Borel $\sigma$-algebra} is $\BorelAlg{X} = \GenSigmaAlg{\mathcal{T}}$, i.e., the $\sigma$-algebra generated by the open sets. Elements of $\BorelAlg{X}$ are called \textbf{Borel sets}.
\end{definition}

\begin{proposition}
On $\bb{R}$, the Borel $\sigma$-algebra can be generated by:
\begin{enumerate}
    \item Open intervals: $\GenSigmaAlg{\{(a, b) \mid a, b \in \bb{R}\}} = \BorelAlg{\bb{R}}$
    \item Half-open intervals: $\GenSigmaAlg{\{(a, b] \mid a, b \in \bb{R}\}} = \BorelAlg{\bb{R}}$
    \item Rays: $\GenSigmaAlg{\{(a, \infty) \mid a \in \bb{R}\}} = \BorelAlg{\bb{R}}$
\end{enumerate}
\end{proposition}

\begin{proof}
(1) Every open set in $\bb{R}$ is a countable union of open intervals. Thus $\GenSigmaAlg{\text{open intervals}} \supseteq \GenSigmaAlg{\text{open sets}} = \BorelAlg{\bb{R}}$. The reverse is clear since open intervals are open.

(2) $(a, b] = \bigcap_{n \in \bb{N}} (a, b + 1/n)$, so half-open intervals are Borel. Conversely, $(a, b) = \bigcup_{n \in \bb{N}} (a, b - 1/n]$.

(3) $(a, \infty)$ is open, so rays are Borel. Conversely, $(a, b) = (a, \infty) \setminus [b, \infty) = (a, \infty) \cap (-\infty, b)$ where $(-\infty, b) = \bigcup_{n \in \bb{N}} (-n, b)$ and $(-n, b) = (-n, \infty) \setminus [b, \infty)$.
\end{proof}

\bigskip
\section{Set Decomposition Lemmas}

The following lemmas allow us to transform arbitrary well-ordered families of sets into families with special structure (disjoint or monotone) while preserving the union or intersection.

\begin{lemma}[Disjointification of Well-Ordered Families]
Let $\lambda$ be an ordinal and $(A_\alpha)_{\alpha < \lambda}$ be a family of sets indexed by ordinals less than $\lambda$. Define:
\[ B_\alpha = A_\alpha \setminus \bigcup_{\beta < \alpha} A_\beta \]
Then:
\begin{enumerate}
    \item $(B_\alpha)_{\alpha < \lambda}$ are pairwise disjoint
    \item Each $B_\alpha \subseteq A_\alpha$
    \item $\bigcup_{\alpha < \lambda} A_\alpha = \bigcup_{\alpha < \lambda} B_\alpha$
\end{enumerate}
\end{lemma}

\begin{proof}
\textbf{Property 1:} Let $\beta < \alpha < \lambda$. We have $B_\alpha = A_\alpha \setminus \bigcup_{\gamma < \alpha} A_\gamma \subseteq A_\alpha \setminus A_\beta$. Since $B_\beta \subseteq A_\beta$, we have $B_\alpha \cap B_\beta \subseteq (A_\alpha \setminus A_\beta) \cap A_\beta = \emptyset$.

\medskip
\textbf{Property 2:} By definition, $B_\alpha = A_\alpha \setminus \bigcup_{\beta < \alpha} A_\beta \subseteq A_\alpha$.

\medskip
\textbf{Property 3:} ($\supseteq$) Each $B_\alpha \subseteq A_\alpha$, so $\bigcup_{\alpha < \lambda} B_\alpha \subseteq \bigcup_{\alpha < \lambda} A_\alpha$.

($\subseteq$) Let $x \in \bigcup_{\alpha < \lambda} A_\alpha$. Define $\gamma = \min\{\alpha < \lambda \mid x \in A_\alpha\}$. This minimum exists since the ordinals are well-ordered. Then $x \in A_\gamma$ and $x \notin A_\beta$ for all $\beta < \gamma$. Thus $x \in A_\gamma \setminus \bigcup_{\beta < \gamma} A_\beta = B_\gamma \subseteq \bigcup_{\alpha < \lambda} B_\alpha$.
\end{proof}

\begin{lemma}[Monotonization to Increasing Family]
Let $\lambda$ be an ordinal and $(A_\alpha)_{\alpha < \lambda}$ be a family of sets. Define:
\[ C_\alpha = \bigcup_{\beta \leq \alpha} A_\beta \]
Then:
\begin{enumerate}
    \item $(C_\alpha)_{\alpha < \lambda}$ is increasing: $\alpha < \beta \implies C_\alpha \subseteq C_\beta$
    \item Each $A_\alpha \subseteq C_\alpha$
    \item $\bigcup_{\alpha < \lambda} A_\alpha = \bigcup_{\alpha < \lambda} C_\alpha$
\end{enumerate}
\end{lemma}

\begin{proof}
\textbf{Property 1:} If $\alpha < \beta$, then $\{\gamma \mid \gamma \leq \alpha\} \subseteq \{\gamma \mid \gamma \leq \beta\}$, so $C_\alpha = \bigcup_{\gamma \leq \alpha} A_\gamma \subseteq \bigcup_{\gamma \leq \beta} A_\gamma = C_\beta$.

\medskip
\textbf{Property 2:} $A_\alpha \subseteq \bigcup_{\beta \leq \alpha} A_\beta = C_\alpha$.

\medskip
\textbf{Property 3:} ($\subseteq$) Each $A_\alpha \subseteq C_\alpha$, so $\bigcup_{\alpha < \lambda} A_\alpha \subseteq \bigcup_{\alpha < \lambda} C_\alpha$.

($\supseteq$) Each $C_\alpha = \bigcup_{\beta \leq \alpha} A_\beta \subseteq \bigcup_{\alpha < \lambda} A_\alpha$.
\end{proof}

\begin{lemma}[Monotonization to Decreasing Family]
Let $\lambda$ be an ordinal and $(A_\alpha)_{\alpha < \lambda}$ be a family of sets. Define:
\[ D_\alpha = \bigcap_{\beta \leq \alpha} A_\beta \]
Then:
\begin{enumerate}
    \item $(D_\alpha)_{\alpha < \lambda}$ is decreasing: $\alpha < \beta \implies D_\alpha \supseteq D_\beta$
    \item Each $D_\alpha \subseteq A_\alpha$
    \item $\bigcap_{\alpha < \lambda} A_\alpha = \bigcap_{\alpha < \lambda} D_\alpha$
\end{enumerate}
\end{lemma}

\begin{proof}
\textbf{Property 1:} If $\alpha < \beta$, then $\{\gamma \mid \gamma \leq \alpha\} \subseteq \{\gamma \mid \gamma \leq \beta\}$, so $D_\alpha = \bigcap_{\gamma \leq \alpha} A_\gamma \supseteq \bigcap_{\gamma \leq \beta} A_\gamma = D_\beta$.

\medskip
\textbf{Property 2:} $D_\alpha = \bigcap_{\beta \leq \alpha} A_\beta \subseteq A_\alpha$.

\medskip
\textbf{Property 3:} ($\supseteq$) Each $D_\alpha \supseteq \bigcap_{\beta < \lambda} A_\beta$, so $\bigcap_{\alpha < \lambda} D_\alpha \supseteq \bigcap_{\alpha < \lambda} A_\alpha$.

($\subseteq$) For each $\gamma < \lambda$, we have $D_\alpha \subseteq A_\gamma$ whenever $\alpha \geq \gamma$. Thus $\bigcap_{\alpha < \lambda} D_\alpha \subseteq A_\gamma$. Since $\gamma$ is arbitrary, $\bigcap_{\alpha < \lambda} D_\alpha \subseteq \bigcap_{\gamma < \lambda} A_\gamma$.
\end{proof}

\begin{remark}
These decompositions are fundamental in measure theory. For countable families:
\begin{itemize}
    \item \textbf{Disjointification} reduces countable unions to disjoint countable unions, essential for $\sigma$-additivity: $\mu(\bigcup_{n \in \bb{N}} A_n) = \mu(\bigcup_{n \in \bb{N}} B_n) = \sum_{n \in \bb{N}} \mu(B_n)$.
    \item \textbf{Monotonization} allows application of measure continuity: $\mu(\bigcup_{n \in \bb{N}} A_n) = \mu(\bigcup_{n \in \bb{N}} C_n) = \lim_{n \to \infty} \mu(C_n)$.
\end{itemize}
The transfinite versions are used in set-theoretic constructions such as the Borel hierarchy.
\end{remark}

\bigskip
\section{$\lambda$-Systems (Dynkin Systems)}

\begin{definition}[$\lambda$-System (First Formulation)]
A collection $\mathcal{L} \subseteq \Powerset{X}$ is a \textbf{$\lambda$-system} (or \textbf{Dynkin system}) if:
\begin{enumerate}
    \item[(I.1)] $X \in \mathcal{L}$
    \item[(I.2)] If $A, B \in \mathcal{L}$ and $A \subseteq B$, then $B \setminus A \in \mathcal{L}$
    \item[(I.3)] If $(A_n)_{n \in \bb{N}} \subseteq \mathcal{L}$ with $A_n \subseteq A_{n+1}$, then $\bigcup_{n \in \bb{N}} A_n \in \mathcal{L}$
\end{enumerate}
\end{definition}

\begin{definition}[$\lambda$-System (Second Formulation)]
Equivalently, $\mathcal{L} \subseteq \Powerset{X}$ is a $\lambda$-system if:
\begin{enumerate}
    \item[(II.1)] $\emptyset \in \mathcal{L}$
    \item[(II.2)] If $A \in \mathcal{L}$, then $A^c \in \mathcal{L}$
    \item[(II.3)] If $(A_n)_{n \in \bb{N}} \subseteq \mathcal{L}$ are pairwise disjoint, then $\bigcup_{n \in \bb{N}} A_n \in \mathcal{L}$
\end{enumerate}
\end{definition}

\begin{proposition}[Equivalence of Formulations]
The two formulations of $\lambda$-system are equivalent.
\end{proposition}

\begin{proof}
\textbf{(I) $\Rightarrow$ (II):}

\textit{(II.1):} $\emptyset = X \setminus X \in \mathcal{L}$ by (I.2) since $X \in \mathcal{L}$ and $X \subseteq X$.

\textit{(II.2):} $A^c = X \setminus A \in \mathcal{L}$ by (I.2) since $A \subseteq X$.

\textit{(II.3):} Let $(A_n)_{n \in \bb{N}}$ be pairwise disjoint. Define $B_n = \bigcup_{k=1}^n A_k$. Then $B_n \subseteq B_{n+1}$ (increasing). We show each $B_n \in \mathcal{L}$ by induction on $n$:
\begin{itemize}
    \item $B_1 = A_1 \in \mathcal{L}$.
    \item $B_{n+1} = B_n \cup A_{n+1}$. From (II.2), $B_n^c \in \mathcal{L}$. Since $A_{n+1} \subseteq B_n^c$ (disjointness), we have $B_n^c \setminus A_{n+1} = (B_n \cup A_{n+1})^c \in \mathcal{L}$ by (I.2). Thus $B_{n+1} = B_n \cup A_{n+1} \in \mathcal{L}$ by (II.2).
\end{itemize}
By (I.3), $\bigcup_{n \in \bb{N}} A_n = \bigcup_{n \in \bb{N}} B_n \in \mathcal{L}$.

\medskip
\textbf{(II) $\Rightarrow$ (I):}

\textit{(I.1):} $X = \emptyset^c \in \mathcal{L}$ by (II.2).

\textit{(I.2):} Let $A \subseteq B$ with $A, B \in \mathcal{L}$. We have $B \setminus A = B \cap A^c$. Consider the disjoint decomposition:
\[ B = A \cup (B \setminus A), \quad B^c, \quad A \setminus B = \emptyset \text{ (since } A \subseteq B) \]
Actually, $B \setminus A$, $A$, and $B^c$ partition $X$. By (II.2), $B^c \in \mathcal{L}$. Now $X = A \cup (B \setminus A) \cup B^c$, i.e., these three are pairwise disjoint. By (II.3), $A \cup (B \setminus A) \cup B^c = X \in \mathcal{L}$. Actually, we need a different approach:

Define $C = B \setminus A$. We must show $C \in \mathcal{L}$. Note $C^c = A \cup B^c$. We have $A$ and $B^c$ are disjoint (since $A \subseteq B$). By (II.3), $A \cup B^c \in \mathcal{L}$. Thus $C^c \in \mathcal{L}$, so $C \in \mathcal{L}$ by (II.2).

\textit{(I.3):} Let $A_n \nearrow A$. By the Disjointification Lemma, define $B_1 = A_1$ and $B_n = A_n \setminus A_{n-1}$ for $n \geq 2$. Then $(B_n)_{n \in \bb{N}}$ are pairwise disjoint and $\bigcup_{n \in \bb{N}} B_n = A$. Each $B_n \in \mathcal{L}$ by (I.2) (which we just proved). By (II.3), $A \in \mathcal{L}$.
\end{proof}

\begin{intuition}
A $\lambda$-system is weaker than a $\sigma$-algebra:
\begin{itemize}
    \item Formulation I: closure under \emph{proper} differences (when $A \subseteq B$) and \emph{increasing} countable unions.
    \item Formulation II: closure under complements and \emph{disjoint} countable unions.
\end{itemize}
Neither requires closure under arbitrary unions or intersections. The key insight: a $\lambda$-system that is \emph{also} a $\pi$-system must be a $\sigma$-algebra. This is the content of Dynkin's $\pi$-$\lambda$ theorem.
\end{intuition}

\begin{proposition}
Every $\sigma$-algebra is a $\lambda$-system.
\end{proposition}

\begin{proof}
Let $\mathcal{M}$ be a $\sigma$-algebra.

\textbf{Axiom 1:} $X = \emptyset^c \in \mathcal{M}$.

\medskip
\textbf{Axiom 2:} $B \setminus A = B \cap A^c \in \mathcal{M}$.

\medskip
\textbf{Axiom 3:} Countable unions are in $\mathcal{M}$ by definition.
\end{proof}

\begin{lemma}[Intersection of $\lambda$-Systems]
Let $\{\mathcal{L}_i\}_{i \in I}$ be an arbitrary family of $\lambda$-systems on $X$. Then $\bigcap_{i \in I} \mathcal{L}_i$ is a $\lambda$-system.
\end{lemma}

\begin{proof}
Let $\mathcal{L} = \bigcap_{i \in I} \mathcal{L}_i$.

\textbf{Axiom 1:} $X \in \mathcal{L}_i$ for all $i$, so $X \in \mathcal{L}$.

\medskip
\textbf{Axiom 2:} Let $A, B \in \mathcal{L}$ with $A \subseteq B$. Then $A, B \in \mathcal{L}_i$ for all $i$. Since each $\mathcal{L}_i$ is a $\lambda$-system, $B \setminus A \in \mathcal{L}_i$ for all $i$. Thus $B \setminus A \in \mathcal{L}$.

\medskip
\textbf{Axiom 3:} Let $(A_n)_{n \in \bb{N}} \subseteq \mathcal{L}$ with $A_n \nearrow$. Then each $A_n \in \mathcal{L}_i$ for all $i$. Since each $\mathcal{L}_i$ is a $\lambda$-system, $\bigcup_{n \in \bb{N}} A_n \in \mathcal{L}_i$ for all $i$. Thus $\bigcup_{n \in \bb{N}} A_n \in \mathcal{L}$.
\end{proof}

\begin{definition}[Generated $\lambda$-System]
For $\mathcal{F} \subseteq \Powerset{X}$, the \textbf{$\lambda$-system generated by $\mathcal{F}$}, denoted $\GenLambda{\mathcal{F}}$, is:
\[ \GenLambda{\mathcal{F}} = \bigcap \{ \mathcal{L} \mid \mathcal{L} \text{ is a } \lambda\text{-system on } X \text{ and } \mathcal{F} \subseteq \mathcal{L} \} \]
\end{definition}

\begin{proposition}[Properties of Generated $\lambda$-Systems]
Let $\mathcal{F} \subseteq \Powerset{X}$. Then:
\begin{enumerate}
    \item $\GenLambda{\mathcal{F}}$ is a $\lambda$-system (by the Intersection Lemma)
    \item $\mathcal{F} \subseteq \GenLambda{\mathcal{F}}$ (extension)
    \item If $\mathcal{L}$ is a $\lambda$-system and $\mathcal{F} \subseteq \mathcal{L}$, then $\GenLambda{\mathcal{F}} \subseteq \mathcal{L}$ (minimality)
\end{enumerate}
\end{proposition}

\begin{proof}
\textbf{Property 1:} By the Intersection Lemma. The intersection is nonempty since $\Powerset{X}$ is a $\lambda$-system containing $\mathcal{F}$.

\medskip
\textbf{Property 2:} Every $\mathcal{L}$ in the intersection satisfies $\mathcal{F} \subseteq \mathcal{L}$. Thus $\mathcal{F} \subseteq \bigcap \mathcal{L} = \GenLambda{\mathcal{F}}$.

\medskip
\textbf{Property 3:} If $\mathcal{L} \supseteq \mathcal{F}$ is a $\lambda$-system, then $\mathcal{L}$ appears in the intersection. Hence $\GenLambda{\mathcal{F}} \subseteq \mathcal{L}$.
\end{proof}

\begin{corollary}[Monotonicity and Idempotence]
Let $\mathcal{F}, \mathcal{G} \subseteq \Powerset{X}$. Then:
\begin{enumerate}
    \item If $\mathcal{F} \subseteq \mathcal{G}$, then $\GenLambda{\mathcal{F}} \subseteq \GenLambda{\mathcal{G}}$ (monotonicity)
    \item $\GenLambda{\GenLambda{\mathcal{F}}} = \GenLambda{\mathcal{F}}$ (idempotence)
\end{enumerate}
\end{corollary}

\begin{proof}
\textbf{Monotonicity:} Let $\mathcal{F} \subseteq \mathcal{G}$. Then $\GenLambda{\mathcal{G}}$ is a $\lambda$-system containing $\mathcal{G} \supseteq \mathcal{F}$. By minimality, $\GenLambda{\mathcal{F}} \subseteq \GenLambda{\mathcal{G}}$.

\medskip
\textbf{Idempotence:} $\GenLambda{\mathcal{F}}$ is a $\lambda$-system, so by minimality, $\GenLambda{\GenLambda{\mathcal{F}}} \subseteq \GenLambda{\mathcal{F}}$. By extension, $\GenLambda{\mathcal{F}} \subseteq \GenLambda{\GenLambda{\mathcal{F}}}$. Thus equality.
\end{proof}

\begin{theorem}[$\lambda$-System Induction Principle]
Let $\mathcal{F} \subseteq \Powerset{X}$ and let $P$ be a property of subsets of $X$. Suppose:
\begin{enumerate}
    \item $P(X)$
    \item $\Forall :{F \in \mathcal{F}}{P(F)}$
    \item $P(A) \text{ and } P(B) \text{ and } A \subseteq B \implies P(B \setminus A)$
    \item $\Forall :{n \in \bb{N}}{P(A_n)} \text{ and } A_n \nearrow A \implies P(A)$
\end{enumerate}
Then $\Forall :{A \in \GenLambda{\mathcal{F}}}{P(A)}$.
\end{theorem}

\begin{proof}
Let $\mathcal{M} = \{A \subseteq X \mid P(A)\}$. By hypotheses (1), (3), (4), $\mathcal{M}$ is a $\lambda$-system. By hypothesis (2), $\mathcal{F} \subseteq \mathcal{M}$. By minimality, $\GenLambda{\mathcal{F}} \subseteq \mathcal{M}$.
\end{proof}

\begin{proposition}[Constructive Description of $\GenLambda{\mathcal{F}}$]
Define $\mathcal{F}_0 = \mathcal{F} \cup \{X\}$ and for $n \geq 0$:
\[ \mathcal{F}_{n+1} = \mathcal{F}_n \cup \{ B \setminus A \mid A, B \in \mathcal{F}_n, A \subseteq B \} \cup \left\{ \bigcup_{k \in \bb{N}} A_k \mid (A_k)_{k \in \bb{N}} \subseteq \mathcal{F}_n, A_k \nearrow \right\} \]
Then $\GenLambda{\mathcal{F}} = \bigcup_{n \in \bb{N}} \mathcal{F}_n$.
\end{proposition}

\begin{proof}
Let $\mathcal{G} = \bigcup_{n \in \bb{N}} \mathcal{F}_n$.

\textbf{$\mathcal{G}$ is a $\lambda$-system:}

\textit{Axiom 1:} $X \in \mathcal{F}_0 \subseteq \mathcal{G}$.

\textit{Axiom 2:} Let $A, B \in \mathcal{G}$ with $A \subseteq B$. Then $A \in \mathcal{F}_m$ and $B \in \mathcal{F}_n$ for some $m, n$. Let $k = \max(m, n)$. Then $A, B \in \mathcal{F}_k$, so $B \setminus A \in \mathcal{F}_{k+1} \subseteq \mathcal{G}$.

\textit{Axiom 3:} Let $(A_k)_{k \in \bb{N}} \subseteq \mathcal{G}$ with $A_k \nearrow$. Each $A_k \in \mathcal{F}_{n_k}$ for some $n_k$. If all $n_k$ are bounded by some $m$, then all $A_k \in \mathcal{F}_m$, so $\bigcup_k A_k \in \mathcal{F}_{m+1} \subseteq \mathcal{G}$.

\medskip
\textbf{$\mathcal{F} \subseteq \mathcal{G}$:} Since $\mathcal{F} \subseteq \mathcal{F}_0 \subseteq \mathcal{G}$.

\medskip
\textbf{$\mathcal{G}$ is minimal:} Let $\mathcal{L}$ be any $\lambda$-system with $\mathcal{F} \subseteq \mathcal{L}$. We show $\mathcal{F}_n \subseteq \mathcal{L}$ for all $n$ by induction.

\textit{Base:} $\mathcal{F}_0 = \mathcal{F} \cup \{X\} \subseteq \mathcal{L}$ since $\mathcal{F} \subseteq \mathcal{L}$ and $X \in \mathcal{L}$.

\textit{Step:} Assume $\mathcal{F}_n \subseteq \mathcal{L}$. If $A, B \in \mathcal{F}_n \subseteq \mathcal{L}$ with $A \subseteq B$, then $B \setminus A \in \mathcal{L}$. If $(A_k) \subseteq \mathcal{F}_n \subseteq \mathcal{L}$ with $A_k \nearrow$, then $\bigcup_k A_k \in \mathcal{L}$. Thus $\mathcal{F}_{n+1} \subseteq \mathcal{L}$.

Hence $\mathcal{G} = \bigcup_{n \in \bb{N}} \mathcal{F}_n \subseteq \mathcal{L}$, proving minimality.

By the three properties, $\mathcal{G} = \GenLambda{\mathcal{F}}$.
\end{proof}

\begin{theorem}[Cardinality of Generated $\lambda$-Systems]
For any $\mathcal{F} \subseteq \Powerset{X}$:
\[ |\GenLambda{\mathcal{F}}| \leq \max(|\mathcal{F}|, \aleph_0)^{\aleph_0} \]
\end{theorem}

\begin{proof}
Let $\kappa = \max(|\mathcal{F}|, \aleph_0)$. We show $|\mathcal{F}_n| \leq \kappa^{\aleph_0}$ for all $n$ by induction.

\textbf{Base:} $|\mathcal{F}_0| = |\mathcal{F} \cup \{X\}| \leq \kappa + 1 = \kappa \leq \kappa^{\aleph_0}$.

\medskip
\textbf{Step:} Assume $|\mathcal{F}_n| \leq \kappa^{\aleph_0}$. Elements of $\mathcal{F}_{n+1} \setminus \mathcal{F}_n$ are either:
\begin{itemize}
    \item Proper differences $B \setminus A$: at most $|\mathcal{F}_n|^2 \leq (\kappa^{\aleph_0})^2 = \kappa^{\aleph_0}$ such elements.
    \item Increasing unions of sequences from $\mathcal{F}_n$: at most $|\mathcal{F}_n|^{\aleph_0} \leq (\kappa^{\aleph_0})^{\aleph_0} = \kappa^{\aleph_0}$ such elements.
\end{itemize}
Thus $|\mathcal{F}_{n+1}| \leq \kappa^{\aleph_0}$.

\medskip
\textbf{Conclusion:} $|\GenLambda{\mathcal{F}}| = |\bigcup_{n \in \bb{N}} \mathcal{F}_n| \leq \aleph_0 \cdot \kappa^{\aleph_0} = \kappa^{\aleph_0}$.
\end{proof}


\bigskip
\section{Monotone Class Lemma}

\begin{theorem}[Monotone Class Lemma]
If $\mathcal{A}$ is an algebra on $X$, then $\GenSigmaAlg{\mathcal{A}} = \GenMonotone{\mathcal{A}}$.
\end{theorem}

\begin{intuition}
The Monotone Class Lemma says:
\begin{center}
\emph{Starting from an algebra, monotone limits give you everything.}
\end{center}

Why is this useful? When proving two $\sigma$-algebras are equal, it sometimes suffices to check that a certain property holds for an algebra and is preserved under monotone limits. This is particularly natural when dealing with measure continuity, where monotone sequences arise naturally.

The proof strategy: show the generated monotone class has enough structure (complementation + unions) to be a $\sigma$-algebra, then minimality forces equality.
\end{intuition}

\begin{proof}
We proceed in three steps. Our goal is to show $\GenMonotone{\mathcal{A}} = \GenSigmaAlg{\mathcal{A}}$.

\textbf{Step 1: One inclusion is immediate.}

Every $\sigma$-algebra is a monotone class (as shown earlier). Thus $\GenSigmaAlg{\mathcal{A}}$ is a monotone class containing $\mathcal{A}$. Since $\GenMonotone{\mathcal{A}}$ is the \emph{smallest} monotone class containing $\mathcal{A}$:
\[ \GenMonotone{\mathcal{A}} \subseteq \GenSigmaAlg{\mathcal{A}} \]

\textbf{Step 2: Showing $\GenMonotone{\mathcal{A}}$ is an algebra.}

For the reverse inclusion, we show $\GenMonotone{\mathcal{A}}$ is a $\sigma$-algebra. We first prove it is an algebra.

\textit{Step 2a: Closure under complements.}

Define:
\[ \mathcal{G} = \{A \in \GenMonotone{\mathcal{A}} \mid A^c \in \GenMonotone{\mathcal{A}}\} \]

We show $\mathcal{G}$ is a monotone class.

\textbf{Increasing sequences:} Let $A_n \nearrow A$ with each $A_n \in \mathcal{G}$. Then:
\begin{itemize}
    \item Each $A_n^c \in \GenMonotone{\mathcal{A}}$ (since $A_n \in \mathcal{G}$)
    \item $A_n^c \searrow A^c$ (taking complements reverses the direction)
    \item Since $\GenMonotone{\mathcal{A}}$ is closed under decreasing limits, $A^c \in \GenMonotone{\mathcal{A}}$
\end{itemize}
Thus $A \in \mathcal{G}$.

\textbf{Decreasing sequences:} Let $A_n \searrow A$ with each $A_n \in \mathcal{G}$. Then:
\begin{itemize}
    \item Each $A_n^c \in \GenMonotone{\mathcal{A}}$
    \item $A_n^c \nearrow A^c$
    \item Since $\GenMonotone{\mathcal{A}}$ is closed under increasing limits, $A^c \in \GenMonotone{\mathcal{A}}$
\end{itemize}
Thus $A \in \mathcal{G}$.

Since $\mathcal{A}$ is an algebra, it is closed under complements, so $\mathcal{A} \subseteq \mathcal{G}$. Since $\mathcal{G}$ is a monotone class containing $\mathcal{A}$:
\[ \GenMonotone{\mathcal{A}} \subseteq \mathcal{G} \subseteq \GenMonotone{\mathcal{A}} \]
Therefore $\GenMonotone{\mathcal{A}} = \mathcal{G}$, proving closure under complements.

\textit{Step 2b: Closure under finite unions.}

We use a double bootstrap similar to the $\pi$-$\lambda$ theorem.

\textbf{First bootstrap:} For fixed $B \in \mathcal{A}$, define:
\[ \mathcal{H}_B = \{A \in \GenMonotone{\mathcal{A}} \mid A \cup B \in \GenMonotone{\mathcal{A}}\} \]

$\mathcal{H}_B$ is a monotone class because:
\begin{itemize}
    \item If $A_n \nearrow A$ with $A_n \in \mathcal{H}_B$, then $A_n \cup B \nearrow A \cup B$. Since each $A_n \cup B \in \GenMonotone{\mathcal{A}}$ and $\GenMonotone{\mathcal{A}}$ is closed under increasing limits, $A \cup B \in \GenMonotone{\mathcal{A}}$.
    \item Similarly for decreasing sequences.
\end{itemize}

Since $\mathcal{A}$ is an algebra, $A \cup B \in \mathcal{A} \subseteq \GenMonotone{\mathcal{A}}$ for any $A \in \mathcal{A}$. Thus $\mathcal{A} \subseteq \mathcal{H}_B$, so $\GenMonotone{\mathcal{A}} = \mathcal{H}_B$.

\textbf{Second bootstrap:} For fixed $A \in \GenMonotone{\mathcal{A}}$, define:
\[ \mathcal{H}_A = \{B \in \GenMonotone{\mathcal{A}} \mid A \cup B \in \GenMonotone{\mathcal{A}}\} \]

$\mathcal{H}_A$ is a monotone class by the same reasoning. From the first bootstrap, for any $B \in \mathcal{A}$, we have $A \cup B \in \GenMonotone{\mathcal{A}}$ (since $A \in \GenMonotone{\mathcal{A}} = \mathcal{H}_B$). Thus $\mathcal{A} \subseteq \mathcal{H}_A$, so $\GenMonotone{\mathcal{A}} = \mathcal{H}_A$.

\textit{Conclusion of Step 2:} For any $A, B \in \GenMonotone{\mathcal{A}}$, we have $A \cup B \in \GenMonotone{\mathcal{A}}$. Together with closure under complements and the fact that $\emptyset \in \mathcal{A} \subseteq \GenMonotone{\mathcal{A}}$, this shows $\GenMonotone{\mathcal{A}}$ is an algebra.

\textbf{Step 3: An algebra that is also a monotone class is a $\sigma$-algebra.}

Let $(A_n)_{n \in \bb{N}} \subseteq \GenMonotone{\mathcal{A}}$. Define:
\[ B_n = \bigcup_{i=1}^n A_i \]

Since $\GenMonotone{\mathcal{A}}$ is an algebra, each $B_n \in \GenMonotone{\mathcal{A}}$ (finite unions). The sequence $(B_n)$ is increasing: $B_n \nearrow \bigcup_{n \in \bb{N}} A_n$.

Since $\GenMonotone{\mathcal{A}}$ is a monotone class (closed under increasing limits):
\[ \bigcup_{n \in \bb{N}} A_n \in \GenMonotone{\mathcal{A}} \]

Thus $\GenMonotone{\mathcal{A}}$ is closed under countable unions. Combined with closure under complements and $\emptyset \in \GenMonotone{\mathcal{A}}$, it is a $\sigma$-algebra.

\textit{Conclusion:} Since $\GenMonotone{\mathcal{A}}$ is a $\sigma$-algebra containing $\mathcal{A}$:
\[ \GenSigmaAlg{\mathcal{A}} \subseteq \GenMonotone{\mathcal{A}} \]
Combined with Step 1: $\GenMonotone{\mathcal{A}} = \GenSigmaAlg{\mathcal{A}}$.
\end{proof}

\bigskip
\section{Dynkin's $\pi$-$\lambda$ Theorem}

\begin{theorem}[Dynkin's $\pi$-$\lambda$ Theorem]
If $\mathcal{P}$ is a $\pi$-system and $\mathcal{L}$ is a $\lambda$-system with $\mathcal{P} \subseteq \mathcal{L}$, then $\GenSigmaAlg{\mathcal{P}} \subseteq \mathcal{L}$.

Equivalently: If $\mathcal{P}$ is a $\pi$-system, then $\GenLambda{\mathcal{P}} = \GenSigmaAlg{\mathcal{P}}$.
\end{theorem}

\begin{intuition}
\textbf{Why is this theorem powerful?} It allows us to prove that two $\sigma$-finite measures agree by checking only a $\pi$-system, which is often much simpler. For example:
\begin{itemize}
    \item To show Lebesgue measure is unique, we only need to check it on intervals.
    \item To prove independence, we only need to check products of generating sets.
\end{itemize}

\textbf{The bootstrap strategy:} The proof uses a clever ``bootstrapping'' argument. We want to show $\GenLambda{\mathcal{P}}$ is closed under intersections. We cannot verify this directly, so we:
\begin{enumerate}
    \item First show: every element of $\GenLambda{\mathcal{P}}$ intersects nicely with elements of $\mathcal{P}$.
    \item Then upgrade: every element of $\GenLambda{\mathcal{P}}$ intersects nicely with \emph{all} elements of $\GenLambda{\mathcal{P}}$.
\end{enumerate}
\end{intuition}

\begin{proof}
We proceed in four steps. Our goal is to show $\GenLambda{\mathcal{P}} = \GenSigmaAlg{\mathcal{P}}$.

\textbf{Step 1: One inclusion is immediate.}

Every $\sigma$-algebra is a $\lambda$-system, so $\GenSigmaAlg{\mathcal{P}}$ is the smallest $\sigma$-algebra containing $\mathcal{P}$, and it is also a $\lambda$-system. Since $\GenLambda{\mathcal{P}}$ is the \emph{smallest} $\lambda$-system containing $\mathcal{P}$, and $\GenSigmaAlg{\mathcal{P}}$ is one such $\lambda$-system:
\[ \GenLambda{\mathcal{P}} \subseteq \GenSigmaAlg{\mathcal{P}} \]

\medskip
\textbf{Step 2: A $\lambda$-system that is also a $\pi$-system is a $\sigma$-algebra.}

We prove this as a lemma within the proof. Let $\mathcal{L}$ be both a $\lambda$-system and a $\pi$-system. We verify the $\sigma$-algebra axioms:

\textit{(1) $\emptyset \in \mathcal{L}$:} Since $\mathcal{L}$ is a $\lambda$-system, $X \in \mathcal{L}$. Thus $\emptyset = X \setminus X \in \mathcal{L}$ (proper difference).

\textit{(2) Closure under complements:} For any $A \in \mathcal{L}$, since $A \subseteq X$, we have $A^c = X \setminus A \in \mathcal{L}$ (proper difference).

\textit{(3) Closure under countable unions:} First, we show $\mathcal{L}$ is closed under finite unions. For $A, B \in \mathcal{L}$:
\[ A \cup B = (A^c \cap B^c)^c \]
Since $A^c, B^c \in \mathcal{L}$ (by closure under complements) and $A^c \cap B^c \in \mathcal{L}$ (by $\pi$-system property), we have $A \cup B \in \mathcal{L}$. By induction, $\mathcal{L}$ is closed under all finite unions.

Now let $(A_n)_{n \in \bb{N}} \subseteq \mathcal{L}$. By the Disjointification Lemma, define $B_1 = A_1$ and $B_n = A_n \setminus \bigcup_{k=1}^{n-1} A_k$ for $n \geq 2$. Then $(B_n)_{n \in \bb{N}}$ are pairwise disjoint and $\bigcup_{n \in \bb{N}} A_n = \bigcup_{n \in \bb{N}} B_n$.

Each $B_n \in \mathcal{L}$: For $n = 1$, $B_1 = A_1 \in \mathcal{L}$. For $n \geq 2$, let $C_n = \bigcup_{k=1}^{n-1} A_k \in \mathcal{L}$ (finite union). Then $A_n \cap C_n \in \mathcal{L}$ (by $\pi$-system property), and $B_n = A_n \setminus (A_n \cap C_n) \in \mathcal{L}$ (proper difference, since $A_n \cap C_n \subseteq A_n$).

By the Monotonization Lemma, $S_n = \bigcup_{k=1}^n B_k$ forms an increasing sequence with $S_n \nearrow \bigcup_{n \in \bb{N}} B_n$. Each $S_n \in \mathcal{L}$ (finite union). By the increasing union property of $\lambda$-systems, $\bigcup_{n \in \bb{N}} A_n = \bigcup_{n \in \bb{N}} B_n \in \mathcal{L}$.

Thus $\mathcal{L}$ is a $\sigma$-algebra. $\square$

\medskip
\textbf{Strategy:} By this lemma, to show $\GenLambda{\mathcal{P}}$ is a $\sigma$-algebra, it suffices to show $\GenLambda{\mathcal{P}}$ is a $\pi$-system, i.e., closed under finite intersections.

\medskip
\textbf{Step 3: The first bootstrap---intersecting with the generating $\pi$-system.}

Define the collection:
\[ \mathcal{G}_1 = \{A \in \GenLambda{\mathcal{P}} \mid \Forall :{P \in \mathcal{P}}{A \cap P \in \GenLambda{\mathcal{P}}}\} \]

We claim $\mathcal{G}_1$ is a $\lambda$-system.

\textit{Verification of $\lambda$-system axioms for $\mathcal{G}_1$:}

\textbf{Axiom 1} ($X \in \mathcal{G}_1$): For any $P \in \mathcal{P}$, we have $X \cap P = P \in \mathcal{P} \subseteq \GenLambda{\mathcal{P}}$. Thus $X \in \mathcal{G}_1$.

\medskip
\textbf{Axiom 2} (Proper differences): Let $A, B \in \mathcal{G}_1$ with $A \subseteq B$. We must show $(B \setminus A) \cap P \in \GenLambda{\mathcal{P}}$ for all $P \in \mathcal{P}$.

Note that $(B \setminus A) \cap P = (B \cap P) \setminus (A \cap P)$. Since $A, B \in \mathcal{G}_1$, both $A \cap P$ and $B \cap P$ are in $\GenLambda{\mathcal{P}}$. Moreover, $A \subseteq B$ implies $A \cap P \subseteq B \cap P$. Since $\GenLambda{\mathcal{P}}$ is a $\lambda$-system, $(B \cap P) \setminus (A \cap P) \in \GenLambda{\mathcal{P}}$.

\medskip
\textbf{Axiom 3} (Increasing unions): Let $A_n \nearrow A$ with each $A_n \in \mathcal{G}_1$. For any $P \in \mathcal{P}$:
\[ A_n \cap P \nearrow A \cap P \]
Since each $A_n \cap P \in \GenLambda{\mathcal{P}}$ and $\GenLambda{\mathcal{P}}$ is closed under increasing unions, we have $A \cap P \in \GenLambda{\mathcal{P}}$. Thus $A \in \mathcal{G}_1$.

\textit{Conclusion of Step 3:} Since $\mathcal{P}$ is a $\pi$-system, for any $P', P \in \mathcal{P}$, we have $P' \cap P \in \mathcal{P} \subseteq \GenLambda{\mathcal{P}}$. Thus $\mathcal{P} \subseteq \mathcal{G}_1$. Since $\mathcal{G}_1 \subseteq \GenLambda{\mathcal{P}}$ by definition and $\mathcal{G}_1$ is a $\lambda$-system containing $\mathcal{P}$:
\[ \GenLambda{\mathcal{P}} \subseteq \mathcal{G}_1 \subseteq \GenLambda{\mathcal{P}} \]
Therefore $\GenLambda{\mathcal{P}} = \mathcal{G}_1$.

\medskip
\textbf{Step 4: The second bootstrap---full closure under intersections.}

Define:
\[ \mathcal{G}_2 = \{A \in \GenLambda{\mathcal{P}} \mid \Forall :{B \in \GenLambda{\mathcal{P}}}{A \cap B \in \GenLambda{\mathcal{P}}}\} \]

By the same argument as Step 3, $\mathcal{G}_2$ is a $\lambda$-system (replace $P \in \mathcal{P}$ with $B \in \GenLambda{\mathcal{P}}$ throughout).

From Step 3, we know $\GenLambda{\mathcal{P}} = \mathcal{G}_1$. This means: for any $P \in \mathcal{P}$ and any $B \in \GenLambda{\mathcal{P}}$, we have $P \cap B \in \GenLambda{\mathcal{P}}$ (since $B \in \mathcal{G}_1$).

This shows $\mathcal{P} \subseteq \mathcal{G}_2$. Since $\mathcal{G}_2$ is a $\lambda$-system:
\[ \GenLambda{\mathcal{P}} \subseteq \mathcal{G}_2 \subseteq \GenLambda{\mathcal{P}} \]
Therefore $\GenLambda{\mathcal{P}} = \mathcal{G}_2$.

\textit{Conclusion:} For all $A, B \in \GenLambda{\mathcal{P}}$, we have $A \cap B \in \GenLambda{\mathcal{P}}$. Thus $\GenLambda{\mathcal{P}}$ is a $\pi$-system. Being both a $\pi$-system and a $\lambda$-system, it is a $\sigma$-algebra.

Since $\GenLambda{\mathcal{P}}$ is a $\sigma$-algebra containing $\mathcal{P}$:
\[ \GenSigmaAlg{\mathcal{P}} \subseteq \GenLambda{\mathcal{P}} \]
Combined with Step 1: $\GenLambda{\mathcal{P}} = \GenSigmaAlg{\mathcal{P}}$.
\end{proof}

\bigskip
\begin{compendium}[Set-Theoretic Foundations]
\textbf{Set Systems Hierarchy:}
\begin{itemize}
    \item \textbf{$\pi$-System:} Closed under finite intersections.
    \item \textbf{Algebra:} Contains $\emptyset$, closed under complements and finite unions.
    \item \textbf{Monotone Class:} Closed under increasing and decreasing countable limits.
    \item \textbf{$\sigma$-Algebra:} Contains $\emptyset$, closed under complements and countable unions.
    \item \textbf{$\lambda$-System:} Contains $X$, closed under proper differences and increasing unions.
\end{itemize}

\textbf{Generated Systems:} For each type, $\mathrm{Gen}(\mathcal{F})$ is the intersection of all systems of that type containing $\mathcal{F}$.

\textbf{Cardinality Bounds:}
\begin{itemize}
    \item $|\GenPi{\mathcal{F}}|, |\GenAlg{\mathcal{F}}| \leq \max(|\mathcal{F}|, \aleph_0)$
    \item $|\GenMonotone{\mathcal{F}}|, |\GenLambda{\mathcal{F}}|, |\GenSigmaAlg{\mathcal{F}}| \leq \max(|\mathcal{F}|, \aleph_0)^{\aleph_0}$
\end{itemize}

\textbf{Key Theorems:}
\begin{itemize}
    \item \textbf{Monotone Class Lemma:} $\GenSigmaAlg{\mathcal{A}} = \GenMonotone{\mathcal{A}}$ for any algebra $\mathcal{A}$.
    \item \textbf{Dynkin $\pi$-$\lambda$:} If $\mathcal{P}$ is a $\pi$-system and $\mathcal{L}$ is a $\lambda$-system with $\mathcal{P} \subseteq \mathcal{L}$, then $\GenSigmaAlg{\mathcal{P}} \subseteq \mathcal{L}$. Equivalently: $\GenLambda{\mathcal{P}} = \GenSigmaAlg{\mathcal{P}}$.
\end{itemize}

\textbf{Set Decomposition:} Any family $(A_\alpha)_{\alpha < \lambda}$ can be transformed into:
\begin{itemize}
    \item \textbf{Disjoint family} $(B_\alpha)$ with same union, where $B_\alpha = A_\alpha \setminus \bigcup_{\beta < \alpha} A_\beta$.
    \item \textbf{Increasing family} $(C_\alpha)$ with same union, where $C_\alpha = \bigcup_{\beta \leq \alpha} A_\beta$.
    \item \textbf{Decreasing family} $(D_\alpha)$ with same intersection, where $D_\alpha = \bigcap_{\beta \leq \alpha} A_\beta$.
\end{itemize}
\end{compendium}

